\section{Related Work}
 
\subsection{Offline RL as sequence modeling}

Offline Reinforcement Learning \cite{offline_rl_survey} is frequently addressed by techniques that are based on value functions \cite{kumar2020conservative,kostrikov2021offline}, goal-conditioned behavior cloning \cite{nair2018visual,ding2019goal} and reward-conditioned behavior cloning \cite{kumar2019reward,srivastava2019training}.
As previously explained, recent works show how this problem can be tackled as a sequence model problem. In \cite{emmons2021rvs} RL via supervised learning is analyzed, this work shows that sometimes pure supervised learning is competitive wrt more complex architectures, given careful tuning of the policy capacity. In contrast, other recent works have adopted the use of Transformers \cite{transformers} as sequence prediction models, and they have provided a promising direction for offline RL, which is also flexible enough to be combined with mentioned techniques. 

In particular, Trajectory Transformers (TTs) \cite{trajectory_transformers} and Decision Transformers (DTs) \cite{decision_transformers} follow this idea by using beam-search planning and reward conditioning, respectively. Notably, these architectures are further studied combining techinques based on Dynamic Programming (such as Q-learning) \cite{yamagata2023q}, online fine-tuning \cite{zheng2022online}, few-shot policy generalization \cite{xu2022prompting}, and automatically-generated waypoints \cite{badrinath2023waypoint}. There are also works that does not build directly on TTs or DTs, such as \cite{huang2024diffusion}, where Transformers are used to overcome inefficiencies of Diffusion Models.

\subsection{Safe Reinforcement Learning}

Safe RL focuses on methods for ensuring that learned policies satisfy some given safety constraints. Online RL techniques \cite{achiam2017constrained,tessler2018reward} typically model safety requirements as constraints on the expected cost, and enforce safety through constrained optimization or Lagrangian formulations during policy learning. Recently, also offline Safe RL is gaining popularity, including techniques for both soft constraints \cite{xu2022constraints} and hard constraints \cite{zheng2024safe}. Adaptations of DTs and TTs are also studied in the context of offline Safe RL \cite{liu2023constrained,wang2024safe,guo2024temporal,zhang2023saformer}, confirming the growing interest in these models in recent years. These works share a common formulation grounded in Constrained MDPs, where safety is enforced by conditioning the policy on cost signals and constraint thresholds. Despite their differences, all these approaches operate by modifying the conditioning signals or relabeling trajectory returns to steer generation toward safe behavior. In contrast, our method leaves the trajectory data and its reward/cost structure untouched and instead regularizes the training objective with a differentiable logic loss derived from DFA progression, injecting temporal-logic knowledge directly into the learning process rather than into the data representation.
Recent work \cite{tian2023reinforcement} also addresses RL using TTs and \ltlf constraints, but their approach differs from ours as it relies on reward shaping derived from DFA states to imitate trajectories.
%They relies on reward shaping derived from DFA states and trains the model to imitate trajectories generated with these shaped rewards. In contrast, our method integrates temporal logic into the training objective through logic-based regularization, enabling constraint-aware learning without modifying the reward function.


An alternative line of work enforces safety through shielding \cite{alshiekh2018safe,jansen2018safe,belardinelli2025probabilistic}, where a runtime monitor prevents the agent from executing actions that may lead to unsafe states. 
While shielding provides strong safety guarantees, it typically assumes the availability of an explicit environment model, it is primarly designed for on-policy learning, and it is structurally limited to the "safety fragment" of formal specifications. In contrast, our method is well suited for offline RL, do only require knowledge of the constraints without assuming any other knowledge of the environment, and can be used to impose any LTLf formula as a constraint, also the ones outside the safety fragment. 

\subsection{Temporally-constrained sequence generation}
\begin{comment}
% old section on RL and LTLf removed

Temporal logic formalisms, such as \ltl \cite{LTL} and \ltlf \cite{LTLf}, are often used in RL with non-Markovian rewards to define temporally extended goals and constraints. These allow to recover the Markovian property by conditioning the policy on the state of an automaton constructed from the formula. Well-known examples include Restraining Bolts \cite{de2019foundations,de2020restraining} and Reward Machines \cite{icarte2022reward,camacho2019ltl,toro2018teaching}. \ltl-based methods are also employed for multi-task RL \cite{vaezipoor2021ltl2action,jackermeier2024deepltl}, reward-shaping techniques \cite{hasanbeig2019certified,shao2023sample}, task inference \cite{hasanbeig2024symbolic}, and shielding \cite{varricchione2024pure,corsi2024verification}.
\end{comment}

Recently, there has been increasing interest in constraining autoregressive sequence generation through logical knowledge. Applications span Cyber-Physical Systems~\cite{STLnet}, Business Process Management~\cite{UmiliPMAI24,DiFrancescomarino17}, and natural language generation with Large Language Models (LLMs)~\cite{trident,llm_beam_search_1,llm_beam_search_2,llm_sampling1,montecarlo_llm,pseudosemantic_loss}.
%
Most prior work incorporates constraints at \emph{test time}, guiding suffix generation via constrained beam search~\cite{DiFrancescomarino17,trident,llm_beam_search_1,llm_beam_search_2}, auxiliary tractable models~\cite{llm_aux_mod_1,llm_aux_mod_2}, or conditioned sampling~\cite{llm_sampling1,montecarlo_llm}.
%
In contrast, we enforce constraints at \emph{training time} through an auxiliary logical loss. Only few approaches follow this direction~\cite{STLnet,UmiliPMAI24,pseudosemantic_loss}. Among them,~\cite{UmiliPMAI24,pseudosemantic_loss} estimate the probability of satisfying a formula via pseudoprobabilities or Monte Carlo, while STLnet~\cite{STLnet} relies on a student–teacher scheme that is difficult to apply in discrete domains~\cite{UmiliPMAI24}.

A key limitation of existing methods is that they target \emph{homogeneous} sequences (e.g., text), assuming that logical formulas are directly expressed over the generated tokens. While reasonable for language, this assumption is limiting in RL settings.
Here, we consider trajectories of states, actions, and rewards (as formalized in Eq.~\ref{eq:trace_form}), where constraints are naturally specified over abstract, high-level \emph{symbols} that must be grounded in subsequences representing states or state–action pairs.
%
Our approach builds on~\cite{UmiliPMAI24,mezini2025neuro}, which approximate \ltlf satisfaction via Monte Carlo and use it as a training loss. We extend this framework to Decision and Trajectory Transformers for safe offline RL applications.

\begin{comment}
    

Zambaldi et al. [5]
showed that augmenting transformers with relational reasoning improve performance in combinatorial
environments and Ritter et al. [69] showed iterative self-attention allowed for RL agents to better
utilize episodic memories. Parisotto et al. [4] discussed design decisions for more stable training of
transformers in the high-variance RL setting. Unlike our work, these still use actor-critic algorithms
for optimization, focusing on novelty in architecture.


This allows us to draw upon the simplicity and scalabil
ity of the Transformer architecture, and associated advances in language modeling
such as GPT-x and BERT. In particular, we present Decision Transformer, an
architecture that casts the problem of RL as conditional sequence modeling. Un
like prior approaches to RL that fit value functions or compute policy gradients,
Decision Transformer simply outputs the optimal actions by leveraging a causally
masked Transformer. By conditioning an autoregressive model on the desired
return (reward), past states, and actions, our Decision Transformer model can gen
erate future actions that achieve the desired return. 


reframed offline RL as sequence modeling, training autoregressive transformers over trajectories and performing planning via decoding (e.g., beam search). 

This work connects four areas: offline RL, sequence-model RL (including trajectory- and decision-transformer approaches), safe/constrained RL, and temporal-logic-guided learning/planning. Our focus differs from step-wise cost shaping by emphasizing finite-trace temporal specifications and automata-level evaluation/decision mechanisms in offline transformer policies. We also emphasize reproducible benchmark protocols and standardized reporting for logic-constrained sequence RL. (Citations to be completed in the final version.)
\end{comment}